% ARTEFACTO A1: Matriz Docente de Auditoría EC3 y Diagnóstico RLC
\documentclass[11pt, a4paper, landscape]{article}

\usepackage[T1]{fontenc}
\usepackage[utf8]{inputenc}
\usepackage{tgpagella}
\usepackage[spanish, es-noshorthands]{babel}
\usepackage{amsmath, amssymb, bm}
\usepackage{tabularx}
\usepackage{booktabs}
\usepackage[most]{tcolorbox}
\usepackage[margin=1.5cm]{geometry}
\usepackage{ragged2e}
\usepackage{fancyhdr}
\usepackage{tikz}

\pagestyle{fancy}
\fancyhf{}
\setlength{\headheight}{16pt}
\lhead{\textbf{UCB Sede Tarija} | Ing. Mecatrónica}
\chead{}
\rhead{IMT-121 Circuitos Electrónicos I | Auditoría EC3}
\lfoot{Prof: M.Sc. Ing. Bernardo Quiroga Turdera}
\rfoot{Página \thepage}
\renewcommand{\headrulewidth}{0.4pt}

\definecolor{ucbblue}{RGB}{0, 51, 102}

\begin{document}

\begin{tcolorbox}[colback=ucbblue!5!white, colframe=ucbblue, title=\textbf{Artefacto A1: Matriz de Evidencia y Diagnóstico Operativo (EC3)}]
    \centering \textbf{Registro Estricto de Capturas y Viabilidad Física — Sesión W14S3}
\end{tcolorbox}

\noindent\textit{Instrucciones: Documente la existencia real de capturas instrumentales trazables. Si no hay evidencia, marque "Pendiente". Diferencie claramente un fallo técnico comprobado de una hipótesis de error.}

\vspace{0.3cm}
{\small
\setlength{\tabcolsep}{3pt}
\renewcommand{\arraystretch}{1.5}
\noindent
\begin{tabularx}{\textwidth}{
  >{\bfseries\RaggedRight\hsize=0.6\hsize}X 
  >{\RaggedRight\arraybackslash\hsize=0.8\hsize}X 
  >{\centering\arraybackslash\hsize=0.6\hsize}X 
  >{\centering\arraybackslash\hsize=0.7\hsize}X 
  >{\RaggedRight\arraybackslash\hsize=0.8\hsize}X 
  >{\centering\arraybackslash\hsize=0.6\hsize}X 
  >{\centering\arraybackslash\hsize=0.6\hsize}X 
  >{\centering\arraybackslash\hsize=0.7\hsize}X 
  >{\centering\arraybackslash\hsize=0.5\hsize}X 
  >{\RaggedRight\arraybackslash\hsize=1.2\hsize}X 
  >{\centering\arraybackslash\hsize=0.7\hsize}X 
  >{\RaggedRight\arraybackslash\hsize=1.2\hsize}X}
\toprule
\textbf{Grupo / Nombres} & \textbf{Captura RC ID} & \textbf{$\tau_{\text{med}}$ [ms] o NO MED} & \textbf{RC Estado} & \textbf{Captura RLC ID} & \textbf{$L_{\text{nom}}$ [mH]} & \textbf{DCR med [$\Omega$]} & \textbf{RLC Estado} & \textbf{GO Topol.} & \textbf{Problema: Causa Conf. / Hipótesis} & \textbf{Acceso Pendiente} & \textbf{Acción / Resp / Fecha} \\
\midrule
 & & & & & & & & & & & \\ \addlinespace
 & & & & & & & & & & & \\ \addlinespace
 & & & & & & & & & & & \\ \addlinespace
 & & & & & & & & & & & \\ \addlinespace
 & & & & & & & & & & & \\ 
\bottomrule
\end{tabularx}
}

\vspace{0.4cm}
\begin{tcolorbox}[colback=white, colframe=black, title=\textbf{Checklist Rápido de Causalidad RLC (Para rellenar columna Problema)}]
\begin{itemize}
    \item (a) Cortocircuito de Tierra (CH1/CH2 GND conectados en nodos distintos al retorno común).
    \item (b) Falla/Polaridad del Capacitor (Ej. Uso de $470\,\mu\text{F}$ electrolítico polarizado en circuito oscilatorio).
    \item (c) DCR y Tolerancia Inductiva excediendo límite crítico $R_{\text{crit}} = 2\sqrt{L/C}$ (Sistema sobreamortiguado accidental).
    \item (d) $R_{\text{gen}}$ excesiva (Carga Hi-Z vs $50\,\Omega$) enlenteciendo el sistema.
    \item (e) Sondas 1X/10X atenuando y modificando capacitancia parásita.
    \item (f) Escala / Trigger erróneo (Buscando frecuencia de kHz en una ventana de ms/div inadecuada).
\end{itemize}
\end{tcolorbox}

\end{document}
